% RESULTADO_RECUPERADO: funtor dimensional integral, eq:cZ-internos,
% eq:longitud-ruta y definición de ell_0; composición con artículo III.
% La prueba reunida y su control algebraico son específicos de esta contribución.
\subsection{Velocidad constitutiva y realización cinemática del reloj}
\label{v:sec:enlace-velocidad-constitutiva}

La propagación radiativa utiliza la sección de velocidad producida por el
transporte angular del mismo estado HMT. La identificación se establece
antes de elegir unidades de representación. Sean $x=A_{\mathrm{rad}}$ e
$y=C^*_{\mathrm{rad}}$ las coordenadas angulares generadas en la
sección~\ref{sec:ii-angulos}, con los canales
\eqref{eq:ii-canales}, y sea
$\mathsf R$ la involución autoadjunta de las dos hojas, con proyectores
de rango uno $P_+$ y $P_-$. En el dominio $x>|y|$, el transporte y sus
canales son
\begin{equation}
 T=\exp(-xI-y\mathsf R)=q_+P_++q_-P_-,
 \qquad q_+=e^{-x-y},\quad q_-=e^{-x+y},\quad 0<q_\pm<1.
 \label{v:eq:velocidad-transporte}
\end{equation}
La respuesta de los bloques temporales de longitudes $90$ y $120$ es
\begin{equation}
 \mathsf V=(I-T^{90})(I-T^{120})^{-1}
       =r_+P_++r_-P_-,\qquad
 r_\pm=\frac{1-q_\pm^{90}}{1-q_\pm^{120}}.
 \label{v:eq:velocidad-respuesta}
\end{equation}
La construcción parte de los canales heredados; los símbolos $x,y$ y
$q_\pm$ designan sus coordenadas y no parámetros ajustados a una velocidad
objetivo.

\begin{proposition}[Identidad de las realizaciones constitutiva y cinemática]
\label{v:prop:identidad-velocidad}
En una misma coordenatización dimensional, la sección $c_{\mathrm{int}}$ del reloj
y la sección $c_{\mathrm{vac}}$ de la respuesta constitutiva coinciden:
\begin{equation}
 c_{\mathrm{int}}=c_{\mathrm{vac}}
   =\widehat c\,u_C,
 \qquad
 \widehat c=(r_+r_-)^{-1}
            =\exp(-\mathcal E),
 \label{v:eq:identidad-velocidad}
\end{equation}
donde
\begin{equation}
 \mathcal E=
 \log\!\bigl[(1-q_+^{90})(1-q_-^{90})\bigr]
 -\log\!\bigl[(1-q_+^{120})(1-q_-^{120})\bigr].
 \label{v:eq:velocidad-lector-simetrico}
\end{equation}
La igualdad conserva el intercambio de hojas.
\end{proposition}
\begin{proof}
Los factores que aparecen en los logaritmos son positivos. Por tanto,
\[
 \mathcal E=\log(r_+r_-)
            =\log\det\mathsf V,
 \qquad \exp(-\mathcal E)=(\det\mathsf V)^{-1}.
\]
El determinante se calcula aquí en la realización de dos hojas de rango
uno. La definición cinemática $c_{\mathrm{int}}=\exp(-\mathcal E)u_C$
y la definición constitutiva $c_{\mathrm{vac}}=(r_+r_-)^{-1}u_C$
son, en consecuencia, la misma sección. El intercambio $P_+\leftrightarrow
P_-$ permuta $r_+$ y $r_-$ y conserva su producto. Finalmente,
$\det T=e^{-2x}$ corresponde al transporte elemental, mientras que
$\det\mathsf V=r_+r_-$ corresponde a su respuesta temporal; ambos
determinantes tienen funciones distintas en esta composición.
\end{proof}

\subsubsection{Longitud de trayectoria y duración elemental}

El marco dimensional positivo $(u_S,u_C,u_T,u_Q)$ induce la base de
longitud $u_L=u_Cu_T$. Para una trayectoria enriquecida $\gamma$ de
$n=|\gamma|$ transiciones y canal constante, los lectores aditivos son
\begin{equation}
 T_{\mathrm{int}}(\gamma)=n u_T,
 \qquad L_{\mathrm{int}}(\gamma)=n\widehat c\,u_L.
 \label{v:eq:velocidad-longitud-reloj}
\end{equation}
Su aditividad se refiere al conteo y a la realización de cada transición:
las rutas enriquecidas conservan adicionalmente orientación y memoria.

\begin{corollary}[Normalización de la longitud elemental]
\label{v:cor:longitud-elemental-constitutiva}
Para $n>0$, $L_{\mathrm{int}}(\gamma)/T_{\mathrm{int}}(\gamma)
=c_{\mathrm{vac}}$. En la normalización $t_0=u_T$,
\begin{equation}
 \ell_0=c_{\mathrm{int}}t_0=\widehat c\,u_L.
 \label{v:eq:longitud-elemental-constitutiva}
\end{equation}
En una coordenatización temporal general $t_0=\tau_0u_T$, con $\tau_0>0$,
la expresión correspondiente es $\ell_0=\tau_0\widehat c\,u_L$.
\end{corollary}
\begin{proof}
La división de \eqref{v:eq:velocidad-longitud-reloj} cancela el entero
$n$ y utiliza $u_L/u_T=u_C$. La fórmula elemental se obtiene
multiplicando la sección de velocidad por $t_0$.
\end{proof}

La generalización a un canal dependiente de la arista conserva la fórmula
aditiva $L_{\mathrm{int}}(\gamma)=\sum_{e\in\gamma}\widehat c(e)u_L$.
Su cociente por $nu_T$ es la media de las velocidades de esas aristas.
La realización radiativa homogénea emplea el canal constante declarado
en \eqref{v:eq:velocidad-longitud-reloj}.

\subsubsection{Cambio de coordenatización y normalización adaptada}

\begin{proposition}[Covariancia de la identificación]
\label{v:prop:covariancia-identidad-velocidad}
Sean $u_C'=d u_C$ y $u_T'=\eta u_T$, con $d,\eta>0$.
La misma velocidad y la misma duración elemental tienen coordenadas
\begin{equation}
 \widehat c'=d^{-1}\widehat c,
 \qquad \tau_0'=\eta^{-1}\tau_0,
 \qquad u_L'=d\eta u_L.
 \label{v:eq:velocidad-cambio-carta}
\end{equation}
En particular, si $c_V=\widehat c_Vu_C^V$ y
$u_C^V=d u_C^{III}$, la igualdad de secciones equivale exactamente a
$d\widehat c_V=\widehat c_{III}$.
\end{proposition}
\begin{proof}
Se expresa cada sección en las dos bases. Se obtiene
$\widehat c'u_C'=\widehat c u_C$ y
$\tau_0'u_T'=\tau_0u_T$. Su producto satisface
\[
 \tau_0'\widehat c' u_L'
 =\frac{\tau_0}{\eta}\frac{\widehat c}{d}
       d\eta u_L
 =\tau_0\widehat c u_L=\ell_0.
\]
La última equivalencia resulta de sustituir la relación entre las bases.
\end{proof}

La elección adaptada $d=\widehat c$ da $u_C'=c_{\mathrm{vac}}$ y
$\widehat c'=1$. Esta elección expresa la sección ya construida en una
base adaptada a ella. El coeficiente espectral en la coordenatización interna sigue
siendo $(r_+r_-)^{-1}$; la inversión espectral utiliza esa coordenatización interna.
Por ejemplo, manteniendo las bases de acción y carga, las coordenadas
constitutivas se transforman como
\[
 \widehat\mu'=d\widehat\mu,
 \qquad \widehat\varepsilon'=d\widehat\varepsilon,
 \qquad
 \widehat\mu'\widehat\varepsilon'=(\widehat c')^{-2}.
\]
En la coordenatización adaptada, $\widehat\mu'=r_-/r_+$ y
$\widehat\varepsilon'=r_+/r_-$, pues las coordenadas internas son
$\widehat\mu=r_-^2$ y $\widehat\varepsilon=r_+^2$.
